% GENERATED FILE. Edit canonical lessons, then run npm run generate:books.
% canonical-source-hash: fnv1a64:56c877bf1ff89800
% canonical-lessons: FR-C265-plus-que, FR-C265-depuis-pendant, FR-C265-revision-experience

\chapter{More Than, and for How Long — plus … que, depuis, pendant}
\label{ch:fr-more-than-and-how-long}
\hlchaptermodality{265}

\begin{chapteropening}
\textbf{By the end of this chapter:} \emph{I can compare two things, and say how long something has been going on or lasted.}

Complete FR-C265-revision-experience: compare two ways to travel, and say how long you have lived somewhere and how long you lived somewhere else.
\end{chapteropening}

\section[plus … que, moins … que]{plus … que — more than}

\label{lesson:FR-C265-plus-que}



\begin{quote}
Say \textquotedblleft{}fast\textquotedblright{}. (\emph{Rapide.}) Say \textquotedblleft{}expensive\textquotedblright{}. (\emph{Cher.})

An experience is easier to describe beside another one. This lesson sets two things side by side.
\end{quote}

\subsection*{You'll want to know: plus … que, moins … que}
To compare two things, put \textbf{plus}, \textquotedblleft{}more\textquotedblright{}, in front of the adjective and \textbf{que}, \textquotedblleft{}than\textquotedblright{}, in front of the second thing:

\begin{quote}
Le train est \textbf{plus} rapide \textbf{que} la voiture. — The train is faster than the car. L'avion est \textbf{plus} cher \textbf{que} le train. — The plane is dearer than the train. Tu es \textbf{plus} jeune \textbf{que} moi. — You are younger than me.
\end{quote}

English often builds the comparison into the adjective, \emph{faster}, \emph{younger}; French builds it with \emph{plus}. After \emph{que} a person takes the strong form you know: \emph{que moi}, \emph{que toi}, \emph{que lui}.

Its opposite is \textbf{moins}, \textquotedblleft{}less\textquotedblright{}: \emph{Le vélo est \textbf{moins} cher \textbf{que} la voiture.} And for \textquotedblleft{}as … as\textquotedblright{}, \textbf{aussi … que}: \emph{Le train est \textbf{aussi} cher \textbf{que} l'avion.}

\begin{grammarlens}[title={Grammar Lens: le plus, and the one that refuses}]
For the most of all, give \emph{plus} the article as well:

\begin{quote}
\emph{le train \textbf{le plus} rapide} — the fastest train \emph{la voiture \textbf{la plus} chère} — the dearest car
\end{quote}

These adjectives stay after the noun, as in \emph{le vin rouge}, and the adjective takes its own article: \emph{le train le plus rapide}, the article twice.

One of the commonest adjectives refuses \emph{plus}. \textbf{Bon} becomes \textbf{meilleur}, \textquotedblleft{}better\textquotedblright{}: \emph{Le café est \textbf{meilleur} que le thé.} Its adverb, \emph{bien}, becomes \textbf{mieux}: \emph{Tu parles \textbf{mieux} que moi.} — You speak better than me.
\end{grammarlens}

\begin{cousinweb}
\textbf{plus} is inherited from Latin \emph{plūs}, \textquotedblleft{}more\textquotedblright{} (EtymDB 2.1), which English kept whole as \textbf{plus}. \textbf{moins} is inherited, through Old French \emph{meins}, from Latin \emph{minus}, \textquotedblleft{}less\textquotedblright{} (EtymDB 2.1), English \textbf{minus}. \textbf{meilleur} is inherited, through Old French \emph{meillor}, from Latin \emph{melior}, \textquotedblleft{}better\textquotedblright{} (EtymDB 2.1), which English kept in \textbf{ameliorate}, \textquotedblleft{}to make better\textquotedblright{}; and \textbf{mieux} goes back to its neighbour \emph{melius}.
\end{cousinweb}

\subsection*{Your turn}
Say these aloud:

\begin{itemize}
\raggedright
  \item \emph{Le train est plus rapide que la voiture.}
  \item the other way round — \emph{La voiture est moins rapide que le train.}
  \item \emph{le train le plus rapide}
  \item a view with a comparison — \emph{Je pense que le vélo est meilleur que la voiture.}
\end{itemize}

\begin{tcolorbox}[breakable,colback=teal!4,colframe=teal!35!black,title={Before you move on}]
Say \textquotedblleft{}the plane is dearer than the train\textquotedblright{}. (\emph{\textbf{L'avion est plus cher que le train.}}) How do you say \textquotedblleft{}the fastest train\textquotedblright{}? (\emph{\textbf{Le train le plus rapide.}}) What does \emph{bon} become? (\emph{\textbf{Meilleur.}})
\end{tcolorbox}
\section[depuis, pendant]{depuis, pendant — how long, and still going}

\label{lesson:FR-C265-depuis-pendant}



\begin{quote}
Say \textquotedblleft{}the train is faster than the car\textquotedblright{}. (\emph{Le train est plus rapide que la voiture.}) Then \textquotedblleft{}I live in Paris\textquotedblright{}. (\emph{J'habite à Paris.})

An experience has a length, and French asks whether it is over.
\end{quote}

\subsection*{You'll want to know: depuis, and the present}
\begin{quote}
\textbf{J'habite} ici \textbf{depuis} trois ans. — I have been living here for three years. \textbf{J'habite} ici \textbf{depuis} janvier. — I have lived here since January.
\end{quote}

Listen to the verb. English says \emph{have been living}; French says \textbf{j'habite}, the plain \textbf{present}, because the living is still going on. \textbf{Depuis} takes a length, \textquotedblleft{}for three years\textquotedblright{}, or a starting point, \textquotedblleft{}since January\textquotedblright{}, with the same present. The question is \emph{depuis quand}, since when, or \emph{depuis combien de temps}, for how long.

Everyday speech everywhere has a second way to say it, with the length in front: \emph{\textbf{Ça fait trois ans que} j'habite ici.}

\begin{grammarlens}[title={Grammar Lens: pendant, and il y a}]
When the stretch has \textbf{ended}, the verb goes into the passé composé and the length takes \textbf{pendant}, \textquotedblleft{}for, during\textquotedblright{}:

\begin{quote}
\textbf{J'ai habité} à Paris \textbf{pendant} deux ans. — I lived in Paris for two years. \emph{(and then I left)}
\end{quote}

One English word, \emph{for}, splits in two. \emph{Depuis} with the present: still going on. \emph{Pendant} with the passé composé: over.

And the \textbf{il y a} you know as \textquotedblleft{}there is\textquotedblright{} also means \textquotedblleft{}ago\textquotedblright{}, with an event:

\begin{quote}
\emph{Je suis arrivé ici \textbf{il y a} trois ans.} — I arrived here three years ago.
\end{quote}

So one life can carry all three: \emph{Je suis arrivé il y a trois ans. J'habite ici depuis trois ans. Avant, j'ai habité à Paris pendant deux ans.}
\end{grammarlens}

\begin{cousinweb}
\textbf{depuis} is \emph{de}, \textquotedblleft{}from\textquotedblright{}, plus \textbf{puis}, \textquotedblleft{}then\textquotedblright{} (EtymDB 2.1): \textquotedblleft{}from then on\textquotedblright{}. \emph{Puis} goes back to a spoken Latin \emph{postius}, built on \emph{post}, \textquotedblleft{}after\textquotedblright{}. \textbf{pendant} is the participle of \textbf{pendre}, \textquotedblleft{}to hang\textquotedblright{}, from Latin \emph{pendēre} (EtymDB 2.1): \emph{pendant deux ans} is, at root, \textquotedblleft{}two years hanging\textquotedblright{}, the stretch while the matter hung open.
\end{cousinweb}

\subsection*{Your turn}
Say these aloud:

\begin{itemize}
\raggedright
  \item \emph{J'habite ici depuis trois ans.}
  \item \emph{Ça fait trois ans que j'habite ici.}
  \item the stretch that is over — \emph{J'ai habité à Paris pendant deux ans.}
  \item \textquotedblleft{}ago\textquotedblright{} — \emph{Je suis arrivé il y a trois mois.}
\end{itemize}

\begin{tcolorbox}[breakable,colback=teal!4,colframe=teal!35!black,title={Before you move on}]
Which tense goes with \emph{depuis}, when it is still going on? (The \textbf{present}.) Which word takes over when the stretch is finished, and with which tense? (\emph{\textbf{Pendant}}, with the \textbf{passé composé}.) What does \emph{il y a} mean with an event? (\textbf{Ago.})
\end{tcolorbox}
\section[(review)]{Review — describing an experience}

\label{lesson:FR-C265-revision-experience}



\begin{quote}
Say \textquotedblleft{}the coffee is better than the tea\textquotedblright{}. (\emph{Le café est meilleur que le thé.}) Then \textquotedblleft{}I have been living here for three years\textquotedblright{}. (\emph{J'habite ici depuis trois ans.})
\end{quote}

\begin{grammarlens}[title={What you've built}]
Two chapters, and an experience can be told whole: how it felt, what you hope next, how it compares, and how long it lasted. \emph{Plus … que}, \emph{moins … que}, \emph{aussi … que}, with \emph{meilleur} for the one adjective that refuses. \emph{Depuis} with the present for what is still going on, \emph{pendant} with the passé composé for what is over, and \emph{il y a} for \textquotedblleft{}ago\textquotedblright{}.

Here is a life in four lines:

\begin{quote}
\emph{J'ai habité à Paris pendant deux ans. C'était une ville plus chère que Lyon, mais j'étais content. J'habite à Lyon depuis janvier. J'espère rester ici.}
\end{quote}

\emph{Rester}, \textquotedblleft{}to stay\textquotedblright{}, is one of the standing verb's children, Latin \emph{re-stāre}: to stay back.
\end{grammarlens}

\subsection*{Your turn}
Say these aloud:

\begin{itemize}
\raggedright
  \item \emph{J'ai habité à Paris pendant deux ans.}
  \item \emph{C'était une ville plus chère que Lyon.}
  \item \emph{J'habite à Lyon depuis janvier.}
  \item and the dream beyond it — \emph{Mon rêve, c'est de vivre près de la mer.}
  \item your own — where you live, since when, and one place you lived before, with \emph{pendant}
\end{itemize}

\begin{tcolorbox}[breakable,colback=teal!4,colframe=teal!35!black,title={Before you move on}]
In \emph{j'habite à Lyon depuis janvier}, why the present? (The living is \textbf{still going on}.) Say \textquotedblleft{}less expensive than the car\textquotedblright{}. (\emph{\textbf{Moins cher que la voiture.}}) And which word says \textquotedblleft{}ago\textquotedblright{}? (\emph{\textbf{Il y a.}})
\end{tcolorbox}
